% فصل چهارم: پیاده‌سازی و ارزیابی نتایج
\chapter{پیاده‌سازی و ارزیابی نتایج}
\label{ch:implementation}

در این فصل، بخش عملی پایان‌نامه تشریح می‌شود. پس از معرفی محیط و ابزار
پیاده‌سازی، حل‌گر مرجعی که برای تولید جواب دقیق معادله برگرز توسعه داده
شده معرفی و راستی‌آزمایی می‌گردد. سپس نتایج حل مستقیم معادله با مدل
زمان‌پیوسته و نتایج کشف پارامترها در مسئله معکوس گزارش می‌شود و در پایان،
یافته‌ها تحلیل و با نتایج مقاله مرجع مقایسه می‌گردند.

\section{محیط و ابزار پیاده‌سازی}
\label{sec:env}

همه کدهای این پژوهش با زبان پایتون نسخه $3.11$ نوشته شده‌اند. برای پیاده‌سازی
شبکه عصبی و مشتق‌گیری خودکار از کتابخانه جکس\LTRfootnote{JAX} استفاده شده است
\cite{jax2018} که امکان ترکیب تبدیل‌های تابعی مانند گرادیان‌گیری، برداری‌سازی
و کامپایل به‌موقع را فراهم می‌کند و برای پژوهش‌های محاسبات علمی گزینه‌ای
سبک و مناسب است. محاسبات عددی کمکی با نامپای و سای‌پای و ترسیم نمودارها با
مت‌پلات‌لیب انجام شده است. بهینه‌سازی دومرحله‌ای شامل پیاده‌سازی دستی
الگوریتم آدام و فراخوانی پیاده‌سازی \lr{L-BFGS-B} از کتابخانه سای‌پای است.

آموزش‌ها روی پردازنده مرکزی (بدون کارت گرافیک) اجرا شده‌اند تا شرایط
بازتولید برای خواننده با امکانات معمولی نیز فراهم باشد. برای تکرارپذیری
نتایج، بذر همه مولدهای اعداد تصادفی ثابت نگه داشته شده و نسخه دقیق کد اجرا‌شده
به‌همراه راهنمای نصب و اجرا در پیوست الف و پوشه \lr{code} ارائه شده است.

\section{حل‌گر مرجع و راستی‌آزمایی آن}
\label{sec:reference}

برای آن‌که دقت شبکه آگاه از فیزیک به‌صورت مستقل سنجیده شود، یک حل‌گر مرجع
برای معادله برگرز پیاده‌سازی شد که هیچ ارتباطی با شبکه عصبی ندارد. در این
حل‌گر، گسسته‌سازی مکانی با روش طیفی فوریه روی ۱۰۲۴ مود و با قانون پالایش
دو ـ سوم برای حذف اثر هم‌پوشانی انجام می‌شود \cite{trefethen2000} و
گام‌زنی زمانی با روش \lr{ETDRK4} (رانگ ـ کوتای مرتبه چهار نمایی) با گام
$dt=0.001$ صورت می‌گیرد که برای جمله پخش سفت، پایداری مناسبی دارد
\cite{kassam2005}. خروجی این حل‌گر روی شبکه $101 \times 1024$ زمانی ـ مکانی
ذخیره و به‌عنوان جواب دقیق در ارزیابی‌ها استفاده شده است.

برای اطمینان از صحت حل‌گر مرجع، دو آزمون مستقل انجام شد. نخست، آزمون
همگرایی: با نصف کردن گام زمانی، تغییر نسبی جواب فقط
$6.2\times 10^{-9}$ و با دو برابر کردن تفکیک مکانی، تغییر نسبی
$4.1\times 10^{-4}$ به دست آمد که رفتار مورد انتظار یک روش مرتبه چهار
زمانی و طیفی مکانی است. دوم، مقایسه با جواب نیمه‌تحلیلی کول ـ هوپف
\cite{hopf1950} که با انتگرال‌گیری عددی مستقیم (بدون هیچ گام‌زنی زمانی)
محاسبه شد: بیشترین اختلاف در ۱۹ نقطه نمونه در لحظه $t=1$ کمتر از
$6\times 10^{-4}$ بود و در نقاط دور از جبهه موج به مرتبه $10^{-8}$ می‌رسید.
این دو آزمون نشان می‌دهد که جواب مرجع با دقتی بسیار بالاتر از دقت مورد
انتظار شبکه عصبی محاسبه شده و مبنای قابل‌اعتمادی برای مقایسه است. گفتنی
است که برای معادله برگرز، جواب دقیق تحلیلی نیز به‌صورت سری فوریه در مراجع
کلاسیک موجود است \cite{benton1972}.

\section{حل مستقیم معادله برگرز}
\label{sec:forward}

مسئله مستقیم به این صورت تعریف شد: با معلوم فرض کردن ضرایب معادله برگرز
($\lambda_1=1$ و $\lambda_2=0.01/\pi$)، جواب $u(t,x)$ در دامنه
$x \in [-1,1]$ و $t \in [0,1]$ با شرط اولیه $u(0,x)=-\sin(\pi x)$ و شرایط
مرزی $u(t,-1)=u(t,1)=0$ به دست آید. شبکه مورد استفاده، یک پرسپترون چندلایه
با ۴ لایه پنهان و ۵۰ نورون در هر لایه با فعال‌ساز تانژانت هیپربولیک است که
در مجموع ۷۸۵۱ پارامتر دارد. داده آموزشی شامل ۲۰۰ نقطه روی شرط اولیه (۱۰۰
نقطه) و مرزها (۱۰۰ نقطه) است و ۱۰۰۰۰ نقطه هم‌مکانی نیز در نظر گرفته شد که
نیمی از آن‌ها به‌صورت یکنواخت از کل دامنه و نیم دیگر به‌صورت متمرکز حول
$x=0$ (محل تشکیل جبهه تند موج) نمونه‌برداری شدند. توزیع نقاط آموزشی در شکل
\ref{fig:fwdpoints} نشان داده شده است. خلاصه تنظیمات آموزش در جدول
\ref{tab:fwdconfig} آمده است.

\begin{figure}[t]
\centering
\includegraphics[width=0.85\textwidth]{forward_points}
\caption{توزیع نقاط آموزشی در مسئله مستقیم: نقاط داده روی شرط اولیه و مرزها و نقاط هم‌مکانی در دامنه}
\label{fig:fwdpoints}
\end{figure}

\begin{table}[t]
\caption{تنظیمات آموزش مدل زمان‌پیوسته در مسئله مستقیم}
\label{tab:fwdconfig}
\centering
\singlespacing
\begin{tabular}{|c|c|}
\hline
\textbf{مؤلفه} & \textbf{مشخصات} \\
\hline
معماری شبکه & ۴ لایه پنهان، ۵۰ نورون در هر لایه، فعال‌ساز تانژانت هیپربولیک \\
\hline
تعداد پارامترها & ۷۸۵۱ \\
\hline
نقاط داده ($N_u$) & ۲۰۰ نقطه (۱۰۰ شرط اولیه و ۱۰۰ شرط مرزی) \\
\hline
نقاط هم‌مکانی ($N_f$) & ۱۰۰۰۰ نقطه (نیمی یکنواخت و نیمی متمرکز حول جبهه) \\
\hline
مقداردهی اولیه & روش گزاویه \\
\hline
بهینه‌سازی & ۶۰۰۰ گام آدام و سپس \lr{L-BFGS} \\
\hline
دقت محاسبات & ممیز شناور ۶۴ بیتی \\
\hline
\end{tabular}
\end{table}

آموزش در دو مرحله انجام شد: ابتدا ۶۰۰۰ گام با بهینه‌ساز آدام و نرخ یادگیری
$10^{-3}$ و سپس ادامه آموزش با روش \lr{L-BFGS} تا سقف ۶۰۰۰ ارزیابی تابع.
مقدار تابع هزینه در پایان مرحله آدام \FwdLossAdam\ و در پایان مرحله
\lr{L-BFGS} \FwdLossLBFGS\ به دست آمد که سهم جملات داده و فیزیک در آن
به‌ترتیب \FwdMSEu\ و \FwdMSEf\ بود. نمودار تاریخچه آموزش در شکل
\ref{fig:loss} رسم شده است که کاهش سریع خطا در مرحله آدام و بهبود نهایی در
مرحله \lr{L-BFGS} را نشان می‌دهد.

\begin{figure}[t]
\centering
\includegraphics[width=0.85\textwidth]{loss_history}
\caption{نمودار تاریخچه تابع هزینه در مسئله مستقیم: کاهش سریع در مرحله آدام و بهبود نهایی با روش شبه‌نیوتنی}
\label{fig:loss}
\end{figure}

برای سنجش دقت، خروجی شبکه روی کل شبکه مرجع $101 \times 1024$ ارزیابی و با
جواب دقیق مقایسه شد. خطای نسبی در نرم $L_2$ برابر \FwdRelErr\ به دست آمد.
مقایسه جواب دقیق و پیش‌بینی شبکه در سه لحظه زمانی $t=0.25$، $0.5$ و $0.75$
در شکل \ref{fig:snapshots} نشان داده شده است. همان‌طور که دیده می‌شود،
شبکه شکل کلی جواب و محل جبهه موج را به‌خوبی بازسازی کرده است. نقشه خطای
مطلق در شکل \ref{fig:errmap} نیز نشان می‌دهد که خطا تقریباً در همه دامنه
ناچیز است و فقط در لایه بسیار باریک جبهه موج (حوالی $x=0$ برای $t>0.3$)
متمرکز شده است؛ ناحیه‌ای که گرادیان جواب در آن بسیار تند است و حتی برای
روش‌های عددی کلاسیک نیز دشوارترین بخش مسئله محسوب می‌شود.

\begin{figure}[t]
\centering
\includegraphics[width=\textwidth]{burgers_snapshots}
\caption{مقایسه جواب دقیق و پیش‌بینی شبکه آگاه از فیزیک در سه لحظه زمانی در مسئله مستقیم معادله برگرز}
\label{fig:snapshots}
\end{figure}

\begin{figure}[t]
\centering
\includegraphics[width=0.85\textwidth]{error_heatmap}
\caption{نقشه خطای مطلق بین جواب دقیق و پیش‌بینی شبکه در دامنه زمانی ـ مکانی؛ خطا عمدتاً در لایه باریک جبهه موج متمرکز است}
\label{fig:errmap}
\end{figure}

در جدول \ref{tab:comparison} نتایج این پایان‌نامه با نتایج گزارش‌شده در
مقاله مرجع \cite{raissi2019} مقایسه شده است. همان‌طور که مشاهده می‌شود،
دقت به دست آمده در این بازتولید کمتر از عدد گزارش‌شده در مقاله است که با
توجه به تفاوت امکانات (آموزش روی پردازنده مرکزی در برابر کارت گرافیک، و
بودجه محدودتر تکرارهای بهینه‌سازی) و تفاوت معماری شبکه، امری مورد انتظار
است. نکته مهم، بازتولید کیفی رفتار روش است: دستیابی به جوابی قابل‌قبول در
کل دامنه، تنها با ۲۰۰ نقطه داده اولیه و مرزی.

\begin{table}[t]
\caption{مقایسه نتیجه مسئله مستقیم با مقاله مرجع}
\label{tab:comparison}
\centering
\singlespacing
\begin{tabular}{|c|c|c|c|c|}
\hline
\textbf{منبع} & \textbf{معماری شبکه} & \textbf{تعداد داده} & \textbf{نقاط هم‌مکانی} & \textbf{خطای نسبی $L_2$} \\
\hline
مقاله مرجع \cite{raissi2019} & ۹ لایه، ۲۰ نورون & ۱۰۰ & ۱۰۰۰۰ & $6.7\times 10^{-4}$ \\
\hline
این پایان‌نامه & ۴ لایه، ۵۰ نورون & ۲۰۰ & ۱۰۰۰۰ & \FwdRelErr \\
\hline
\end{tabular}
\end{table}

\section{کشف پارامترهای معادله برگرز}
\label{sec:inverse-exp}

در مسئله معکوس، ضرایب $\lambda_1$ و $\lambda_2$ در باقیمانده
$f = u_t + \lambda_1 u u_x - \lambda_2 u_{xx}$ مجهول فرض شدند و هدف،
شناسایی آن‌ها از روی مشاهده‌های پراکنده جواب بود. داده آموزشی شامل ۵۰۰۰
نقطه پراکنده در دامنه زمانی ـ مکانی است که نیمی یکنواخت و نیمی متمرکز حول
$x=0$ انتخاب شدند (شکل \ref{fig:invpoints})، زیرا جمله شامل $\lambda_2$
(پخش) عملاً فقط در ناحیه جبهه اثر قابل‌توجه دارد و بدون نمونه‌برداری کافی
از آن ناحیه، شناسایی این ضریب ممکن نیست. شبکه مورد استفاده ۵ لایه پنهان با
۳۰ نورون در هر لایه دارد و پارامترهای $\lambda_1$ و $\lambda_2$ از مقدار
اولیه صفر شروع به آموزش کردند. آموزش مانند مسئله مستقیم دومرحله‌ای (۶۰۰۰
گام آدام و سپس \lr{L-BFGS}) انجام شد.

\begin{figure}[t]
\centering
\includegraphics[width=0.85\textwidth]{inverse_points}
\caption{توزیع ۵۰۰۰ نقطه مشاهده پراکنده مورد استفاده در مسئله معکوس؛ رنگ نقاط نشان‌دهنده مقدار جواب است}
\label{fig:invpoints}
\end{figure}

دو سناریو آزموده شد: داده تمیز (بدون نویز) و داده آلوده به نویز گاوسی یک
درصد. نتایج در جدول \ref{tab:inverse} خلاصه شده است. در حالت بدون نویز،
ضریب $\lambda_1$ برابر \LamIClean\ (با خطای \LamIErrClean) و ضریب
$\lambda_2$ برابر \LamIIClean\ (با خطای \LamIIErrClean) شناسایی شد؛ مقادیر
واقعی به‌ترتیب $1$ و $0.01/\pi \approx 0.00318$ هستند. در حالت نویزی نیز
$\lambda_1$ برابر \LamINoisy\ (خطای \LamIErrNoisy) و $\lambda_2$ برابر
\LamIINoisy\ (خطای \LamIIErrNoisy) به دست آمد.

\begin{table}[t]
\caption{نتایج شناسایی پارامترهای معادله برگرز در مسئله معکوس}
\label{tab:inverse}
\centering
\singlespacing
\begin{tabular}{|c|c|c|c|c|}
\hline
\textbf{سناریو} & \textbf{$\lambda_1$ شناسایی‌شده} & \textbf{خطای $\lambda_1$} & \textbf{$\lambda_2$ شناسایی‌شده} & \textbf{خطای $\lambda_2$} \\
\hline
بدون نویز & \LamIClean & \LamIErrClean & \LamIIClean & \LamIIErrClean \\
\hline
نویز یک درصد & \LamINoisy & \LamIErrNoisy & \LamIINoisy & \LamIIErrNoisy \\
\hline
مقدار واقعی & \multicolumn{2}{|c|}{$1$} & \multicolumn{2}{|c|}{$0.01/\pi \approx 0.00318$} \\
\hline
\end{tabular}
\end{table}

همان‌طور که دیده می‌شود، ضریب انتقال ($\lambda_1$) در هر دو سناریو با دقت
قابل‌قبولی شناسایی شده، اما خطای ضریب پخش ($\lambda_2$) به‌مراتب بیشتر است. این
یافته با گزارش‌های مقاله مرجع نیز سازگار است، زیرا در آنجا نیز خطای
$\lambda_2$ همواره بزرگ‌تر از خطای $\lambda_1$ گزارش شده است. دلیل این
امر، کوچکی مقدار $\lambda_2$ و موضعی بودن اثر آن (محدود به لایه باریک
جبهه) است که شناسایی دقیق آن را ذاتاً به مسئله‌ای حساس تبدیل می‌کند؛ در
بخش بعد بیشتر به این موضوع می‌پردازیم.

\section{بحث و تحلیل نتایج}
\label{sec:discussion}

\textbf{کارایی داده‌ای.} مهم‌ترین نقطه قوت مشاهده‌شده، کارایی داده‌ای روش
است: در مسئله مستقیم، تنها با ۲۰۰ نقطه داده روی شرط اولیه و مرز (در برابر
بیش از یکصدهزار نقطه شبکه مرجع)، جواب کل دامنه بازسازی شد. این ویژگی دقیقاً
همان چیزی است که شبکه‌های آگاه از فیزیک را از مدل‌های صرفاً داده‌محور
متمایز می‌کند و در کاربردهایی که داده تجربی کمیاب و گران است اهمیت زیادی
دارد.

\textbf{چالش جبهه تند.} متمرکز شدن خطا در ناحیه جبهه موج نشان می‌دهد که
گلوگاه اصلی دقت، تفکیک‌پذیری لایه‌های تند جواب است. تمرکز نیمی از نقاط
هم‌مکانی در این ناحیه کمک قابل‌توجهی به بهبود نتیجه کرد، اما همچنان دقیق‌ترین
بخش جواب همان ناحیه باقی ماند. این مشاهده با ادبیات موضوع سازگار است و
راهکارهایی مانند وزن‌دهی تطبیقی جملات تابع هزینه یا نمونه‌برداری تطبیقی
نقاط در پژوهش‌های بعدی پیشنهاد شده‌اند.

\textbf{حساسیت شناسایی ضریب پخش.} خطای بیشتر $\lambda_2$ نسبت به
$\lambda_1$ ریشه فیزیکی دارد: جمله پخش $\lambda_2 u_{xx}$ به‌سبب کوچکی
$\lambda_2$، تقریباً در همه دامنه (به‌جز لایه باریک جبهه) در برابر جمله
انتقال قابل‌چشم‌پوشی است؛ در نتیجه سیگنال آموزشی موجود برای تعیین دقیق آن
ضعیف است. هر خطای کوچکی در بازسازی انحنای جواب در ناحیه جبهه نیز مستقیماً
به خطای $\lambda_2$ تبدیل می‌شود. با این حال، مرتبه بزرگی ضریب به‌درستی
شناسایی شد و مقدار آن در هر دو سناریو در محدوده قابل‌قبولی از مقدار واقعی
قرار گرفت.

\textbf{پایداری در برابر نویز.} مقایسه دو سناریوی تمیز و نویزی نشان می‌دهد
که افزودن یک درصد نویز به داده‌ها، دقت شناسایی را به‌طور چشمگیری تخریب
نکرده است. دلیل این پایداری آن است که مشتق‌های لازم برای تشکیل باقیمانده،
به‌جای تفاضل محدود روی داده‌های نویزی، با مشتق‌گیری خودکار از روی شبکه‌ای
هموار محاسبه می‌شوند و شبکه در عمل نقش یک فیلتر هموارساز را نیز ایفا می‌کند.

\textbf{مقایسه با روش‌های کلاسیک.} در مقایسه با حل‌گر مرجع طیفی که جوابی
بسیار دقیق ولی فقط برای یک شرط اولیه و یک مجموعه پارامتر خاص تولید می‌کند،
شبکه آموزش‌دیده تابعی پیوسته و مشتق‌پذیر از جواب در اختیار می‌گذارد که
ارزیابی آن در هر نقطه دلخواه (حتی خارج از شبکه) ارزان است. از سوی دیگر،
هزینه آموزش شبکه به‌مراتب بیشتر از یک اجرای حل‌گر کلاسیک است و دقت نهایی
آن نیز کمتر است. بنابراین در شرایط فعلی، مزیت اصلی این چارچوب نه جایگزینی
حل‌گرهای کلاسیک در مسائل مستقیم استاندارد، بلکه توانایی آن در مسائل معکوس،
مسائل با داده ناقص و ترکیب داده و مدل فیزیکی است.

\textbf{محدودیت‌ها.} دامنه بخش عملی این پژوهش محدود به معادله برگرز
یک‌بعدی و مدل زمان‌پیوسته بود و به‌سبب محدودیت توان محاسباتی، امکان جست‌وجوی
گسترده ابرپارامترها (معماری شبکه، تعداد نقاط، راهبرد بهینه‌سازی) وجود
نداشت. همچنین برآورد عدم قطعیت، که از مزیت‌های بالقوه روش‌های یادگیری
ماشین است، در این پژوهش بررسی نشد.

\section{جمع‌بندی فصل}
\label{sec:ch4-summary}

در این فصل، پیاده‌سازی مدل زمان‌پیوسته شبکه آگاه از فیزیک برای معادله
برگرز ارائه شد. یک حل‌گر مرجع طیفی ـ نمایی توسعه یافت و با دو آزمون مستقل
راستی‌آزمایی گردید. در مسئله مستقیم، جواب کل دامنه تنها با ۲۰۰ نقطه داده
اولیه و مرزی بازسازی شد و خطای نسبی \FwdRelErr\ به دست آمد. در مسئله
معکوس، ضرایب معادله از روی ۵۰۰۰ مشاهده پراکنده شناسایی شدند و اثر نویز بر
دقت شناسایی بررسی گردید. تحلیل نتایج نشان داد که گلوگاه اصلی، تفکیک لایه
تند جبهه موج و حساسیت ذاتی شناسایی ضریب کوچک پخش است.
